\documentclass[12pt]{article}
\usepackage[utf8]{inputenc}
\usepackage[T1]{fontenc}
\usepackage{hyperref}
\title{Xhosa Language Resources for K-12 Curriculum Development}
\author{Compiled by Claude Sonnet 5.5 (claude-sonnet-5-5) for Andr\'as Kornai}
\date{}
\begin{document}
\maketitle

\section{Introduction}
isiXhosa (ISO 639-3 xho, Glottocode xhos1239) is an official language of South Africa with about 11 million speakers, and it is used as a language of instruction. Resources below are those that appeared in my searches. They are mostly curriculum documents, school dictionaries and Bible translations. My searches did not cover grammars, a national lexicography dictionary, corpora or audio, so this list is far from complete.

\section{Primary Educational Resources}
\begin{itemize}
\item CAPS (Curriculum and Assessment Policy Statement) isiXhosa Home Language, Grades 4--6: \url{https://wcedeportal.westerncape.gov.za/eresource/curriculum-and-assessment-policy-statement-caps-isixhosa-home-language-grade-4-6}
\item CAPS isiXhosa Home Language, Grades 7--9: \url{https://wcedeportal.westerncape.gov.za/eresource/curriculum-and-assessment-policy-statement-caps-isixhosa-home-language-grade-7-9}
\item WCED Teaching and Assessment Plans and lesson plans for isiXhosa Home Language (Grades 1, 5--8, 10--12), for example the Grade 1 lesson plan: \url{https://wcedeportal.co.za/eresource/186881}
\item Oxford Bilingual School Dictionary: IsiXhosa and English, Grade 4--9: \url{https://oxford.co.za/shop/resources-and-home-learning/resource-book/9780195766820-oxford-bilingual-school-dictionary-isixhosa-and-english-grade-4-9}
\item Oxford First Bilingual Dictionary: IsiXhosa and English, 2nd edition (ePDF), for the early grades: \url{https://oxford.co.za/?p=60614}
\item Oxford First Bilingual Dictionary flash cards, isiXhosa: \url{https://oxford.co.za/shop/resources-and-home-learning/resource-book/9780190401375-oxford-first-bilingual-dictionary-flash-cards-full-pack-isixhosa/}
\end{itemize}

\section{Linguistic and Reference Materials}
\begin{itemize}
\item Oxford English--Xhosa Dictionary, with a guide to noun classes: \url{https://oxford.co.za/shop/resources-and-home-learning/resource-book/9780199043422-oxford-english-xhosa-dictionary/}
\end{itemize}

\section{Religious and Community Resources}
\begin{itemize}
\item The Bible in isiXhosa, history of the translation (Bible Society of South Africa): \url{https://biblesa.co.za/nso/methopo/about-bible-translation/bible-translations/the-bible-in-isixhosa-2}
\item Bible Society of South Africa, isiXhosa Bible resource page: \url{https://www.biblesociety.co.za/index.php/resources/the-bible?id=117}
\item IsiXhosa Bible 1996 translation on YouVersion: \url{https://www.bible.com/my-MM/versions/3526}
\end{itemize}

\section{Digital and Audio Resources}
I did not find audio, video, corpus or software resources in my searches.

\section{Curriculum Development Context}
South Africa's CAPS provides a national isiXhosa Home Language curriculum. The documents above come from Western Cape Education Department portals, not the national department site, so the official current versions should be checked. Published school dictionaries cover the early and intermediate grades. Gaps in what I found include senior-phase and FET textbooks, grammar references, corpora and audio, which need further searching.

\end{document}